\documentclass[11pt]{article}
\usepackage[margin=1.1in]{geometry}
\usepackage{amsmath,amssymb,amsthm}
\usepackage[colorlinks=true,linkcolor=blue,citecolor=blue]{hyperref}

\newtheorem{theorem}{Theorem}[section]
\newtheorem{lemma}[theorem]{Lemma}
\newtheorem{proposition}[theorem]{Proposition}
\newtheorem{corollary}[theorem]{Corollary}
\newtheorem{definition}[theorem]{Definition}
\newtheorem{remark}[theorem]{Remark}
\newtheorem{problem}[theorem]{Problem}
\newtheorem{conjecture}[theorem]{Conjecture}

\newcommand{\D}{\mathbb{D}}
\newcommand{\T}{\mathbb{T}}
\newcommand{\E}{\mathbb{E}}
\newcommand{\PP}{\mathbb{P}}
\newcommand{\wind}{\operatorname{wind}}

\title{A parity obstruction to the stability of\\random banded Toeplitz matrices}
\author{Yitzchak Shmalo\\
\small Einstein Institute of Mathematics, The Hebrew University of Jerusalem, Jerusalem, Israel\\
\small }
\date{June 2026}

\begin{document}
\maketitle

\begin{abstract}
Let $T_n(P_{r,s})$ denote the $n\times n$ banded Toeplitz matrices with random
symbol $P_{r,s}(z)=\sum_{j=-r}^{s}\xi_jz^{j}$, where the $\xi_j$ are i.i.d.\ and
$m:=r+s$.  By the classical Gohberg--Feldman/B\"ottcher--Silbermann criterion, the
sequence $(T_n)$ is \emph{stable} --- invertible with uniformly bounded inverses
--- exactly when the symbol has no zeros on the unit circle and winding number
zero, i.e.\ when the polynomial $G_m(z)=z^{r}P_{r,s}(z)$ has exactly $r$ zeros in
the open unit disk.  A recent preprint of Manrique
[arXiv:2602.12581] asserts that for balanced bands ($r=s$) stability holds with
probability close to one, while unbalanced bands are exponentially unstable.  We
prove an unconditional \emph{parity obstruction}: for any i.i.d.\ coefficients
with zero mean, unit variance and finite third moment --- including the
Rademacher case --- the number of zeros of $G_m$ in the disk is congruent mod $2$
to the indicator $\{G_m(1)G_m(-1)<0\}$, which is asymptotically a fair coin.
At odd degrees this endpoint sign has an exact finite-degree formula: it is
fair up to the atom where the even- and odd-coefficient sums have equal absolute
value, and is exactly fair for non-atomic laws.  For Littlewood coefficients the
finite-degree endpoint sign is given by central-binomial expressions in both
parities; for Gaussian coefficients one can compute the finite-degree endpoint
bias at every degree.  These endpoint quantities also give finite-degree,
parity-specific caps on every point probability \(\PP(\nu=k)\) and on every
band stability event, while reversal gives an exact finite-degree equality
between complementary band splits.
Consequently
\[
\PP\bigl(\,(T_n(P_{r,s}))_n\ \text{is stable}\,\bigr)\;\le\;\tfrac12+o(1)
\qquad(m\to\infty),
\]
\emph{for every choice of} $(r,s)$.  In particular the balanced-band statement of
[arXiv:2602.12581, Theorem~2.1] cannot hold once its constant $c_3$ is below
$1/2$, and we locate the gap in its proof: the zero-count estimate proved there
confines the number of disk zeros to a window of width $\asymp m^{9/10}$ around
$m/2$, which contains $\asymp m^{9/10}$ integers and cannot pin the count to the
single value $r$.  On the positive side we show, assembling the estimates that
are correct in [arXiv:2602.12581] with the operator-theoretic instability bounds,
that strongly unbalanced bands ($|r-s|>2m^{9/10}$) are exponentially unstable with
probability $1-O((\log m)^{-1/2})$.  The true asymptotics of the stability
probability is governed by a local limit theorem for the number of zeros of a Kac
polynomial in the unit disk, which we state as an open problem.
\end{abstract}

\section{Introduction}\label{sec:intro}

A banded Toeplitz matrix is constant along its diagonals and supported on
finitely many of them: given nonnegative integers $r,s$ and complex numbers
$x_{-r},\dots,x_s$, the $n\times n$ matrix $T_n$ has entries
$(T_n)_{jk}=x_{k-j}$ for $-r\le k-j\le s$ and $0$ otherwise.  The associated
\emph{symbol} is the Laurent polynomial $P(z)=\sum_{j=-r}^{s}x_jz^j$, and the
asymptotic spectral theory of $(T_n(P))_{n\ge1}$ is governed by $P$; the
definitive reference is B\"ottcher--Grudsky \cite{BG}.  A sequence $(A_n)$ of
$n\times n$ matrices is called \emph{stable} if there are $n_0$ and
$M<\infty$ with $A_n$ invertible and $\|A_n^{-1}\|\le M$ for all $n\ge n_0$;
stability of $(T_n(P))$ is equivalent to the invertibility of the associated
Toeplitz operator and holds if and only if
\begin{equation}\label{eq:stability}
P(t)\ne0\ \ \text{for all }t\in\T,
\qquad\text{and}\qquad
\wind(P)=0,
\end{equation}
where $\T$ is the unit circle and $\wind(P)$ the winding number of
$P|_\T$ about the origin \cite[Chapter~3]{BG}.  Writing
\[
G_m(z):=z^{r}P(z)=\sum_{j=0}^{m}x_{j-r}z^{j},
\qquad m:=r+s,
\]
one has $\wind(P)=\nu-r$ whenever $P$ does not vanish on $\T$, where
\[
\nu\;:=\;\#\{z\in\D:\ G_m(z)=0\}
\]
(zeros counted with multiplicity, $\D$ the open unit disk).  Thus
\eqref{eq:stability} reads: \emph{no zeros on the circle, and exactly $r$ of the
$m$ zeros inside the disk}.  Moreover, when $\wind(P)=\omega\ne0$ and $P\ne0$ on
$\T$, the inverses blow up exponentially:
$\|T_n(P)^{-1}\|\ge C_Pe^{\alpha n}$ for every
$\alpha<\min\{\log(1/\delta),\log\mu\}$, where $\delta<1<\mu$ are the moduli of
the zeros of $G_m$ nearest to $\T$ from inside and outside
\cite[Theorem~4.1]{BG}.

It is natural to randomize: let the coefficients $\xi_{-r},\dots,\xi_s$ be
i.i.d.\ random variables, so that $G_m$ is a Kac polynomial of degree $m$.  The
zeros of Kac polynomials concentrate near $\T$ as $m\to\infty$
\cite{IZ,HN}, and for centered coefficients roughly half of them lie inside the
disk.  The stability question then becomes a delicate question about the
\emph{exact} number of zeros in $\D$.

This question was raised, and an answer proposed, in the recent preprint of
Manrique \cite{Manrique}.  Its Theorem~2.1 asserts a dichotomy: for balanced
bands ($r=s$) the sequence $(T_n)$ is stable --- with smallest singular value
bounded below uniformly in $n$ --- with probability $1-c_3-O((\log m)^{-1/2})$,
where $c_3\in(0,1)$ ``can be small''; while for $r\ne s$ the matrices are
exponentially ill-conditioned.  The mechanism proposed there is the zero-count
estimate \cite[Theorem~3.1]{Manrique}: for a broad class $\mathbf H$ of
coefficient laws,
\begin{equation}\label{eq:window}
\PP\Bigl(\Bigl|\frac{\nu}{m}-\frac12\Bigr|\ge m^{-1/10}\Bigr)
=O\Bigl(\frac{(\log m)^2}{m^{1/10}}\Bigr),
\end{equation}
which confines $\nu$ to the window $[\,m/2-m^{9/10},\,m/2+m^{9/10}\,]$ with high
probability.

The present paper makes three points.

\subsection*{1.  A parity obstruction}

Our main result is an unconditional upper bound of one half on the stability
probability, valid for \emph{every} band and a minimal set of moment
assumptions.

\begin{theorem}[Parity obstruction]\label{thm:parity}
Let $\xi_0,\xi_1,\dots,\xi_m$ be i.i.d.\ real random variables with
$\E\xi_0=0$, $\E\xi_0^2=1$ and $\E|\xi_0|^3<\infty$, let
$G_m(z)=\sum_{j=0}^m\xi_jz^j$ and let $\nu$ be the number of zeros of $G_m$ in
$\D$, counted with multiplicity.  Then, as $m\to\infty$:
\begin{enumerate}
\item[(i)] {\rm(deterministic identity)}  On the event $G_m(1)G_m(-1)\ne0$,
\[
\nu\;\equiv\;\mathbf 1\{G_m(1)G_m(-1)<0\}\pmod 2 .
\]
\item[(ii)] {\rm(fair coin)}
$\PP\bigl(G_m(1)G_m(-1)<0\bigr)\longrightarrow\tfrac12$ and
$\PP\bigl(G_m(1)G_m(-1)=0\bigr)=O(m^{-1/2})$.
\item[(iii)] Consequently, for every integer $k=k(m)$,
\[
\PP(\nu=k)\;\le\;\tfrac12+o(1),
\]
and for every choice of nonnegative integers $r=r(m)$, $s=s(m)$ with $r+s=m$,
\[
\PP\bigl(\,(T_n(P_{r,s}))_{n\ge1}\ \text{is stable}\,\bigr)\;\le\;\tfrac12+o(1).
\]
\end{enumerate}
\end{theorem}

The proof occupies one page: zeros of a real polynomial off the real axis come in
conjugate pairs of equal modulus, so the parity of $\nu$ equals the parity of the
number of zeros in the real interval $(-1,1)$, which --- with multiplicity ---
is odd exactly when $G_m(1)$ and $G_m(-1)$ have opposite signs; and
$\bigl(G_m(1),G_m(-1)\bigr)/\sqrt{m+1}$ converges to a pair of \emph{independent}
standard Gaussians, because the two coefficient sums are nearly orthogonal.  No
absolute continuity is required: the result covers the Rademacher (Littlewood)
case.

At odd degrees the endpoint sign has a finite-sample version which does not use
the central limit theorem.

\begin{corollary}[Exact endpoint sign at odd degrees]\label{cor:odd-endpoint}
Let $\xi_0,\dots,\xi_m$ be i.i.d.\ real random variables and let
$m=2q+1$.  Put
\[
E_q:=\sum_{\ell=0}^{q}\xi_{2\ell},
\qquad
O_q:=\sum_{\ell=0}^{q}\xi_{2\ell+1}.
\]
Then
\[
\PP\bigl(G_m(1)G_m(-1)<0\bigr)
=\PP\bigl(|O_q|>|E_q|\bigr)
=\frac12\Bigl(1-\PP\bigl(|E_q|=|O_q|\bigr)\Bigr)
\le \frac12.
\]
If the law of $E_q$ is atomless, in particular if the law of $\xi_0$ is
non-atomic, then
\[
\PP\bigl(G_m(1)G_m(-1)<0\bigr)=\frac12.
\]
\end{corollary}

\begin{proof}
The sums $E_q$ and $O_q$ are independent and identically distributed.  Since
$G_m(1)=E_q+O_q$ and $G_m(-1)=E_q-O_q$,
\[
G_m(1)G_m(-1)=E_q^2-O_q^2,
\]
so the product is negative exactly on $\{|O_q|>|E_q|\}$.  The two events
$\{|O_q|>|E_q|\}$ and $\{|E_q|>|O_q|\}$ have the same probability, and together
with $\{|E_q|=|O_q|\}$ they partition the sample space.  This gives the formula.
If $E_q$ is atomless, then conditioning on $O_q$ gives
$\PP(E_q=O_q)=\PP(E_q=-O_q)=0$, so the equality event has probability zero.  The
non-atomicity of $\xi_0$ implies the non-atomicity of $E_q$ by conditioning on
the remaining summands.
\end{proof}

\begin{corollary}[Exact Littlewood endpoint sign at odd degrees]\label{cor:littlewood-odd}
Assume that the coefficients are independent Rademacher signs and let
\(m=2q+1\).  Put \(n=q+1=(m+1)/2\), and define
\[
  \eta_n:=
  \begin{cases}
  4^{-n}\binom{n}{n/2}^{2},& n\ {\rm even},\\[2mm]
  0,& n\ {\rm odd}.
  \end{cases}
\]
Then
\[
\PP\bigl(G_m(1)G_m(-1)<0\bigr)
=\frac12-\frac{\binom{2n}{n}}{4^n}+\frac{\eta_n}{2}.
\]
In particular,
\[
\PP\bigl(G_m(1)G_m(-1)<0\bigr)
=\frac12-\frac{1}{\sqrt{\pi n}}+O(n^{-1})
=\frac12-\sqrt{\frac{2}{\pi(m+1)}}+O(m^{-1}).
\]
\end{corollary}

\begin{proof}
By Corollary~\ref{cor:odd-endpoint}, it remains to compute
\(\PP(|E_q|=|O_q|)\).  Both \(E_q\) and \(O_q\) have the law of
\[
S_n:=\sum_{i=1}^{n}\varepsilon_i,
\]
where the \(\varepsilon_i\)'s are independent Rademacher signs.  If
\[
p_k:=\PP\{S_n=2k-n\}=2^{-n}\binom nk,\qquad 0\le k\le n,
\]
then, for independent copies \(S_n,S_n'\),
\[
\PP(|S_n|=|S_n'|)
=2\sum_{k=0}^{n}p_k^2-\eta_n .
\]
Indeed, \(|2k-n|=|2\ell-n|\) means either \(\ell=k\) or \(\ell=n-k\), with the
single overlap \(k=\ell=n/2\) only when \(n\) is even.  Vandermonde's identity
gives
\[
\sum_{k=0}^{n}p_k^2
=4^{-n}\sum_{k=0}^{n}\binom nk^2
=4^{-n}\binom{2n}{n}.
\]
Substituting this into
\(\PP(G_m(1)G_m(-1)<0)=\tfrac12(1-\PP(|E_q|=|O_q|))\) gives the exact formula.
The asymptotic form follows from Stirling's formula.
\end{proof}

\begin{corollary}[Exact Littlewood endpoint sign at even degrees]\label{cor:littlewood-even}
Assume that the coefficients are independent Rademacher signs and let
\(m=2q\).  Define
\[
  \theta_q:=
  \begin{cases}
  4^{-q}\binom{q}{q/2}^{2},& q\ {\rm even},\\[2mm]
  0,& q\ {\rm odd}.
  \end{cases}
\]
Then
\[
\PP\bigl(G_m(1)G_m(-1)<0\bigr)=\frac12-\frac{\theta_q}{2}.
\]
Thus the endpoint sign is exactly fair when \(m\equiv2\pmod4\), while for
\(m\equiv0\pmod4\),
\[
\PP\bigl(G_m(1)G_m(-1)<0\bigr)
=\frac12-\frac{2}{\pi m}+O(m^{-2}).
\]
\end{corollary}

\begin{proof}
Write
\[
O_q:=\sum_{\ell=0}^{q-1}\varepsilon_{2\ell+1},\qquad
C_q:=\sum_{\ell=0}^{q-1}\varepsilon_{2\ell},\qquad
E_q:=C_q+\varepsilon_{2q}.
\]
Then \(O_q\) and \(C_q\) are independent copies of a \(q\)-sign sum, and
\(\varepsilon_{2q}\) is independent of both.  Since
\[
G_m(1)G_m(-1)=E_q^2-O_q^2,
\]
the endpoint sign is negative exactly when \(|O_q|>|E_q|\).

Condition on \(O_q\) and \(C_q\).  The two numbers \(|O_q|\) and \(|C_q|\) have
the same parity.  If \(|O_q|>|C_q|\), then
\(|O_q|>|C_q+\varepsilon_{2q}|\) for both choices of the last sign; if
\(|O_q|<|C_q|\), it holds for neither choice.  If \(|O_q|=|C_q|>0\), it holds
for exactly one choice of the last sign.  If \(O_q=C_q=0\), it holds for neither
choice.  Therefore
\[
\PP_{\varepsilon_{2q}}\{|O_q|>|C_q+\varepsilon_{2q}|\}
=\mathbf1_{\{|O_q|>|C_q|\}}
  +\frac12\mathbf1_{\{|O_q|=|C_q|>0\}} .
\]
Averaging and using that \(O_q,C_q\) are independent with the same law gives
\[
\PP\{|O_q|>|E_q|\}
=\frac{1-\PP\{|O_q|=|C_q|\}}2
  +\frac{\PP\{|O_q|=|C_q|\}-\PP\{O_q=C_q=0\}}2
=\frac12-\frac12\PP\{O_q=0\}^2.
\]
The central atom of a \(q\)-sign sum is zero when \(q\) is odd and
\(2^{-q}\binom{q}{q/2}\) when \(q\) is even.  This gives the exact formula.
The asymptotic for \(m=2q\equiv0\pmod4\) follows from Stirling's formula:
\[
\frac12\left(2^{-q}\binom{q}{q/2}\right)^2
=\frac1{\pi q}+O(q^{-2})
=\frac{2}{\pi m}+O(m^{-2}).
\]
\end{proof}

\begin{corollary}[Exact Gaussian endpoint bias]\label{cor:gaussian-endpoint}
Assume that the coefficients are independent standard real Gaussians.  Then
\[
\PP\bigl(G_m(1)G_m(-1)<0\bigr)
=\frac12-\frac1\pi
  \arcsin\!\left(\frac{1}{m+1}\sum_{j=0}^{m}(-1)^j\right).
\]
Equivalently,
\[
\PP\bigl(G_m(1)G_m(-1)<0\bigr)=\frac12
\qquad(m\ \text{odd}),
\]
whereas
\[
\PP\bigl(G_m(1)G_m(-1)<0\bigr)
=\frac12-\frac1\pi\arcsin\!\left(\frac1{m+1}\right)
\qquad(m\ \text{even}).
\]
In the balanced Toeplitz case \(m=2r\), the endpoint obstruction is therefore
slightly biased below a fair coin by
\(\pi^{-1}\arcsin((2r+1)^{-1})=O(r^{-1})\).
\end{corollary}

\begin{proof}
The normalized vector
\[
X:=\frac{G_m(1)}{\sqrt{m+1}},
\qquad
Y:=\frac{G_m(-1)}{\sqrt{m+1}}
\]
is centered bivariate normal with unit variances and correlation
\[
\rho=\E[XY]=\frac{1}{m+1}\sum_{j=0}^{m}(-1)^j .
\]
For a centered bivariate normal pair with correlation \(\rho\), write
\[
X=U,\qquad Y=\rho U+\sqrt{1-\rho^2}\,V,
\]
where \(U,V\) are independent standard Gaussians.  In polar coordinates
\((U,V)=R(\cos\theta,\sin\theta)\), the angle \(\theta\) is uniform on
\([0,2\pi)\).  The two lines \(X=0\) and \(Y=0\) cut the circle into sectors;
the total angular measure of the sectors on which \(XY<0\) is
\(\pi-2\arcsin\rho\).  Hence
\[
\PP(XY<0)=\frac{\pi-2\arcsin\rho}{2\pi}
=\frac12-\frac1\pi\arcsin\rho .
\]
Finally \(\sum_{j=0}^m(-1)^j=0\) for odd \(m\) and \(=1\) for even \(m\).
\end{proof}

\begin{corollary}[Finite-degree parity envelope]\label{cor:finite-envelope}
Let
\[
a_m:=\PP\bigl(G_m(1)G_m(-1)<0\bigr),
\qquad
z_m:=\PP\bigl(G_m(1)G_m(-1)=0\bigr).
\]
Then, for every integer \(k\),
\[
\PP(\nu=k)\le
\begin{cases}
a_m+z_m,& k\ {\rm odd},\\[1mm]
1-a_m,& k\ {\rm even}.
\end{cases}
\]
Moreover, for every band split \(r+s=m\),
\[
\PP\bigl((T_n(P_{r,s}))_{n\ge1}\ {\rm is\ stable}\bigr)
\le
\begin{cases}
a_m,& r\ {\rm odd},\\[1mm]
1-a_m-z_m,& r\ {\rm even}.
\end{cases}
\]
Thus any exact finite-degree formula for the endpoint sign, such as the
Littlewood and Gaussian formulas above, immediately gives a finite-degree
parity cap rather than only the asymptotic \(\tfrac12+o(1)\) bound.
\end{corollary}

\begin{proof}
Let \(Z=\{G_m(1)G_m(-1)=0\}\) and \(A=\{G_m(1)G_m(-1)<0\}\).  On \(Z^c\), the
deterministic parity identity gives
\[
\nu\equiv\mathbf1_A\pmod2 .
\]
Hence \(\{\nu=k\}\subseteq A\cup Z\) when \(k\) is odd, and
\(\{\nu=k\}\subseteq A^c\) when \(k\) is even.  This proves the point-mass
bounds.

For stability, \eqref{eq:stability} and the winding-number formula imply that
the stable event is contained in \(\{\nu=r\}\cap Z^c\).  If \(r\) is odd, the
parity identity forces \(A\).  If \(r\) is even, it forces
\(G_m(1)G_m(-1)>0\), whose probability is \(1-a_m-z_m\).
\end{proof}

\begin{corollary}[Gaussian and Littlewood finite-degree caps]\label{cor:special-envelope}
The envelope in Corollary~\ref{cor:finite-envelope} has the following explicit
forms.
\begin{enumerate}
\item[(i)] For standard real Gaussian coefficients, put
\[
\alpha_m:=\frac1\pi
  \arcsin\!\left(\frac{1}{m+1}\sum_{j=0}^{m}(-1)^j\right).
\]
Then \(z_m=0\), \(a_m=\tfrac12-\alpha_m\), and
\[
\PP(\nu=k)\le
\begin{cases}
\tfrac12-\alpha_m,& k\ {\rm odd},\\[1mm]
\tfrac12+\alpha_m,& k\ {\rm even},
\end{cases}
\qquad
\PP({\rm stable})\le
\begin{cases}
\tfrac12-\alpha_m,& r\ {\rm odd},\\[1mm]
\tfrac12+\alpha_m,& r\ {\rm even}.
\end{cases}
\]
In particular, when \(m\) is odd, every point probability and every band
stability probability is at most \(\tfrac12\) at finite degree.

\item[(ii)] For Littlewood coefficients with \(m=2n-1\) odd, let
\[
\beta_n:=\frac{\binom{2n}{n}}{4^n}-\frac{\eta_n}{2},
\]
where \(\eta_n\) is as in Corollary~\ref{cor:littlewood-odd}.  Then
\[
\PP(\nu=k)\le \frac12+\beta_n\qquad\text{for all }k,
\]
while every band split \(r+s=m\) satisfies the strictly sub-half finite-degree
stability cap
\[
\PP({\rm stable})\le\frac12-\beta_n
=\frac12-\sqrt{\frac{2}{\pi(m+1)}}+O(m^{-1}).
\]

\item[(iii)] For Littlewood coefficients with \(m=2q\) even, let
\[
\gamma_q:=\frac{\theta_q}{2},
\]
where \(\theta_q\) is as in Corollary~\ref{cor:littlewood-even}.  Then \(z_m=0\),
\(a_m=\tfrac12-\gamma_q\), and the Gaussian-form bounds in {\rm(i)} hold with
\(\alpha_m\) replaced by \(\gamma_q\).  Thus all caps are exactly
\(\tfrac12\) when \(q\) is odd, while for even \(q\) the deviation is
\[
\gamma_q=\frac{2}{\pi m}+O(m^{-2}).
\]
\end{enumerate}
\end{corollary}

\begin{proof}
For Gaussian coefficients, \(G_m(1)\) and \(G_m(-1)\) have a non-degenerate
joint density, so \(z_m=0\), and Corollary~\ref{cor:gaussian-endpoint} gives
the displayed value of \(a_m\).  Corollary~\ref{cor:finite-envelope} then gives
the bounds.

For Littlewood coefficients with \(m=2n-1\), Corollary~\ref{cor:littlewood-odd}
gives \(a_m=\tfrac12-\beta_n\).  The endpoint-zero event is
\(\{|E_q|=|O_q|\}\), and the proof of Corollary~\ref{cor:littlewood-odd} gives
its probability \(z_m=2\beta_n\).  Hence the point-mass envelope is
\(\tfrac12+\beta_n\) in both parities, while the stability envelope is
\(\tfrac12-\beta_n\) in both parities.

For Littlewood coefficients with \(m=2q\), the even and odd coefficient sums have
opposite parity, so \(z_m=0\).  Corollary~\ref{cor:littlewood-even} gives
\(a_m=\tfrac12-\gamma_q\), and another application of
Corollary~\ref{cor:finite-envelope} proves the stated bounds.  The asymptotic
forms of \(\beta_n\) and \(\gamma_q\) are those already proved in
Corollaries~\ref{cor:littlewood-odd} and \ref{cor:littlewood-even}.
\end{proof}

Theorem~\ref{thm:parity} contradicts the balanced-band half of
\cite[Theorem~2.1]{Manrique} whenever its constant $c_3$ is taken below
$\tfrac12-o(1)$, as its Remark~2.2 envisages.  Section~\ref{sec:gap} locates the
gap in the argument given there: the step deducing ``$r=m/2$'' from
\eqref{eq:window} treats the window $m\cdot m^{-1/10}=m^{9/10}$ as if it were
$m^{-1/10}$; the window in fact contains $\asymp m^{9/10}$ integers, and
\eqref{eq:window} gives no information about the probability that $\nu$ hits any
single one of them.  Indeed, by Theorem~\ref{thm:parity} that probability cannot
exceed $\tfrac12+o(1)$, no matter how the window shrinks.

\subsection*{2.  Strong asymmetry is correctly unstable}

The unbalanced half of the dichotomy of \cite{Manrique} is correct as soon as the
asymmetry exceeds the width of the window \eqref{eq:window}; the proof is an
assembly of \eqref{eq:window} with the operator-theoretic exponential instability
bound and with the circle-clearance estimate for the zeros of $G_m$ proved in
\cite{KMS} and used in \cite{Manrique}.

\begin{theorem}[Strongly unbalanced bands]\label{thm:unbalanced}
Assume the coefficient law satisfies the hypotheses $\mathbf H$ of
\cite{Manrique} {\rm(}so that \eqref{eq:window} and the circle-clearance estimate
of {\rm\cite[Theorem~1.1]{KMS}} hold{\rm)}.  Let $r+s=m$ with
\[
\bigl|r-s\bigr|\;>\;2\,m^{9/10}.
\]
Then with probability $1-O\bigl((\log m)^{-1/2}\bigr)$ the symbol has no zeros on
$\T$, its winding number is nonzero, and there is a random constant
$C=C(G_m)>0$ such that
\[
\bigl\|T_n(P_{r,s})^{-1}\bigr\|\;\ge\;C\,e^{\alpha_m n}
\qquad\text{for all }n\ge1,
\qquad
\alpha_m:=\log\Bigl(1+\frac{1}{2m^{2}(\log m)^{3}}\Bigr).
\]
In particular $(T_n(P_{r,s}))$ is not stable, and the condition number
$\kappa(T_n)=\sigma_1(T_n)\,\|T_n^{-1}\|$ grows exponentially in $n$.
\end{theorem}

\subsection*{3.  The remaining problem is a local limit theorem}

Theorems \ref{thm:parity} and \ref{thm:unbalanced} leave exactly one regime
genuinely open: bands with $|r-s|=O(m^{9/10})$, and above all the balanced case,
where the stability probability is
\[
p_m\;:=\;\PP(\nu=r)-\PP(\nu=r,\ \exists\,\text{zero on }\T)
\]
and everything hinges on the point probabilities of the zero count $\nu$.  The
parity bound shows $p_m\le\tfrac12+o(1)$; we believe the truth is much smaller.

\begin{conjecture}[Gaussian edge fluctuations]\label{conj:llt}
For standard Gaussian coefficients, there is a constant \(c>0\) such that
\[
\sigma_m^2:=\operatorname{Var}(\nu)\sim c\,m
\]
and
\[
\frac{\nu-m/2}{\sqrt{c\,m}}\;\Longrightarrow\;N(0,1)
\]
with a local limit theorem
\[
\PP(\nu=k)\asymp m^{-1/2}
\]
uniformly for \(|k-m/2|=O(\sqrt m)\).  In particular
\(\max_k\PP(\nu=k)\to0\).  We expect at least this qualitative
anti-concentration conclusion, and hence instability of balanced random banded
Toeplitz matrices with probability \(1-o(1)\), for every coefficient law
satisfying the hypotheses of Theorem~\ref{thm:parity}.
\end{conjecture}

In the Gaussian case the upper half of the conjectured order is a theorem.

\begin{proposition}[Variance upper bound]\label{prop:varupper}
For standard Gaussian coefficients $\xi_j\sim N(0,1)$,
\[
\sigma_m^2=\operatorname{Var}(\nu)=O(m).
\]
\end{proposition}

We prove this in Section~\ref{sec:open} from the Kac--Rice variance identity; the
point is that the boundary shell carries no $|z-w|^{-2}$ singularity in the
\emph{connected} two-point function, which is what makes the real-axis logarithm
of Maslova absent from the disk count.  The matching lower bound
$\sigma_m^2=\Omega(m)$ and the sharp constant remain open
(Problem~\ref{prob:variance}); Monte Carlo puts the constant near $0.30$.

The qualitative conclusion --- balanced bands unstable with probability tending
to one, $\max_k\PP(\nu=k)\to0$ --- would follow from a genuine
anti-concentration theorem, for instance a CLT/LLT or an Esseen-type bound; the
predicted \(m^{-1/2}\) rate for \(p_m\) is the stronger local-limit form.  For
Gaussian coefficients Kac--Rice gives an explicit second-moment
formula,
\[
\operatorname{Var}(\nu)
=\int_{\D}\rho_1+
\int_{\D\times\D}(\rho_2-\rho_1\otimes\rho_1),
\]
which yields the $O(m)$ bound of Proposition~\ref{prop:varupper}; the remaining
difficulty is the matching lower bound and the constant, where one must control
the sign of the connected two-point term at the edge scale, where the unit circle
cuts through the zero cloud.

What is known about the fluctuations of $\nu$?  The expectation is $m/2$ by the
reversal symmetry $\xi_j\mapsto\xi_{m-j}$ (Lemma~\ref{lem:reversal}); the bound
\eqref{eq:window} gives $|\nu-m/2|\le m^{9/10}$ with high probability; and the
parity of $\nu$ is asymptotically a fair coin by Theorem~\ref{thm:parity} ---
which already rules out any concentration of $\nu$ on a single value, and is, to
our knowledge, the only unconditional anti-concentration information available
for $\nu$ at this generality.  Several pieces of evidence point to an edge
fluctuation rather than a bulk fluctuation.  After the rescaling
$z\mapsto1-z/m$, the zeros near \(\T\) are modeled by the Peres--Vir\'ag
determinantal process \cite{PV,HKPV} with Bergman kernel
$\mathbb K(z,w)=\pi^{-1}(1-z\bar w)^{-2}$; for that process the number of points
in $\{|z|<r\}$ has variance $r^2/(1-r^4)$ \cite{Buckley}, which is of order \(m\)
at the cutoff \(r=1-1/m\).  The Shepp--Vanderbei intensity \cite{SV} likewise
identifies an order-\(m\) population of zeros in the \(1/m\)-shell around
\(\T\).  These comparisons support the \(\sigma_m\asymp\sqrt m\) picture, but
they do not by themselves prove it for the disk count of the finite Kac
polynomial.

The logarithmic real-zero fluctuation is a useful warning rather than a substitute
for the disk-count problem.  By Lemma~\ref{lem:parity}, the parity of \(\nu\) is
the parity of the real-zero count in \((-1,1)\), whose variance is
$\bigl(\tfrac4\pi(1-\tfrac2\pi)+o(1)\bigr)\log m$ with a central limit theorem, by
Maslova \cite{Maslova} and Do--Nguyen \cite{DoNguyen}.  This controls the parity
observable, but it does not determine the variance of the two-dimensional disk
count.  Interior rigidity results for related Gaussian analytic functions, such
as bounded variance in annuli away from the boundary \cite{Yeager}, similarly
indicate that the open difficulty is the edge shell.  We record the expected
statement as Conjecture~\ref{conj:llt} and discuss partial routes (Kac--Rice
second moments in the Gaussian case; decorrelation of winding increments along
well-separated arcs) in Section~\ref{sec:open}.

\subsection*{Relation to the literature}

For \emph{dense} (non-banded) random Toeplitz matrices the smallest singular
value and condition number are studied in \cite{condrefs}; for random
\emph{circulant} matrices polynomial conditioning is known \cite{condrefs}.  The
operator-theoretic input \eqref{eq:stability} and the exponential instability
under nonzero winding are classical \cite{BG}.  The zero-count window
\eqref{eq:window} for general coefficient laws is the contribution of
\cite{Manrique}, generalizing the Littlewood case treated in the references
cited there; the circle-clearance estimate is from \cite{KMS}.  Our
Theorem~\ref{thm:parity} appears to be new even for Littlewood polynomials, and
we note that it is of independent interest in the theory of random polynomials:
\emph{the number of zeros of a Kac polynomial in the unit disk is asymptotically
equidistributed mod 2}.

\section{The parity obstruction}\label{sec:parity}

Throughout this section $G(z)=G_m(z)=\sum_{j=0}^m\xi_jz^j$ with
$\xi_j$ i.i.d., $\E\xi_0=0$, $\E\xi_0^2=1$, $\beta:=\E|\xi_0|^3<\infty$.  Note
$G$ has real coefficients, and $G(1)=\sum_j\xi_j$,
$G(-1)=\sum_j(-1)^j\xi_j$.

\begin{lemma}[Parity identity]\label{lem:parity}
Let $g$ be any polynomial with real coefficients and $g(1)g(-1)\ne0$, and let
$\nu(g)$ be the number of zeros of $g$ in $\D$ counted with multiplicity.  Then
\[
\nu(g)\;\equiv\;\mathbf 1\{g(1)g(-1)<0\}\pmod 2 .
\]
\end{lemma}

\begin{proof}
Since $g$ has real coefficients, its non-real zeros come in conjugate pairs
$z,\bar z$ of equal multiplicity and equal modulus; each pair contributes $0$ or
$2$ to $\nu(g)$.  Hence
\[
\nu(g)\;\equiv\;\#\{\text{real zeros of }g\text{ in }(-1,1),
\text{ with multiplicity}\}\pmod2 .
\]
Factor $g(x)=c\prod_i(x-a_i)^{m_i}\cdot q(x)$ over $\mathbb R$, where the $a_i$
are the distinct real zeros and $q$ is a product of irreducible quadratics, so
$q>0$ on $\mathbb R$.  Then
\[
\operatorname{sign}\bigl(g(1)g(-1)\bigr)
=\prod_i\operatorname{sign}\bigl((1-a_i)^{m_i}(-1-a_i)^{m_i}\bigr)
=\prod_{i:\,a_i\in(-1,1)}(-1)^{m_i},
\]
because for $a_i\notin[-1,1]$ the factors $(1-a_i)$ and $(-1-a_i)$ have equal
sign, while for $a_i\in(-1,1)$ they have opposite signs (and $a_i=\pm1$ is
excluded).  Thus $g(1)g(-1)<0$ if and only if
$\sum_{a_i\in(-1,1)}m_i$ is odd.
\end{proof}

\begin{lemma}[Joint CLT for the two boundary values]\label{lem:clt}
As $m\to\infty$,
\[
\frac{1}{\sqrt{m+1}}\bigl(G(1),\,G(-1)\bigr)
\;\Longrightarrow\;(Z_1,Z_2),
\]
where $Z_1,Z_2$ are independent standard Gaussians.  Moreover
$\PP(G(1)=0)\le 4\beta\,m^{-1/2}$ and likewise for $G(-1)$, for all large $m$.
\end{lemma}

\begin{proof}
For fixed $(\alpha,\beta')\in\mathbb R^2$,
$\alpha G(1)+\beta'G(-1)=\sum_{j=0}^m c_j\xi_j$ with
$c_j=\alpha+(-1)^j\beta'$.  The $c_j$ are bounded and
$\sum_jc_j^2=(m+1)(\alpha^2+\beta'^2)+O(1)$, so Lindeberg's condition holds
trivially and
$\bigl(\alpha G(1)+\beta'G(-1)\bigr)/\sqrt{m+1}\Rightarrow
N\bigl(0,\alpha^2+\beta'^2\bigr)$; note
$\operatorname{Cov}(G(1),G(-1))=\sum_{j=0}^m(-1)^j\in\{0,1\}$, which after
normalization vanishes in the limit.  By Cram\'er--Wold the pair converges to
$(Z_1,Z_2)$ as claimed.

For the atom: by the Berry--Esseen theorem, for every $\varepsilon>0$,
\[
\PP\bigl(|G(1)|\le\varepsilon\sqrt{m+1}\bigr)
\le\Phi(\varepsilon)-\Phi(-\varepsilon)+\frac{2C_{\mathrm{BE}}\beta}{\sqrt{m+1}}
\le\varepsilon+\frac{2C_{\mathrm{BE}}\beta}{\sqrt{m+1}},
\]
and letting $\varepsilon\downarrow0$ gives
$\PP(G(1)=0)\le2C_{\mathrm{BE}}\beta(m+1)^{-1/2}\le4\beta m^{-1/2}$ for large
$m$, using $C_{\mathrm{BE}}<1$ \cite{Shevtsova}.
\end{proof}

\begin{proof}[Proof of Theorem \ref{thm:parity}]
(i) is Lemma~\ref{lem:parity} applied to $g=G_m$.  For (ii):
the set $\{(x,y):xy<0\}$ has boundary of Lebesgue measure zero in
$\mathbb R^2$ and $\PP\bigl((Z_1,Z_2)\in\partial\{xy<0\}\bigr)=0$, so by
Lemma~\ref{lem:clt} and the portmanteau theorem,
$\PP(G(1)G(-1)<0)\to\PP(Z_1Z_2<0)=\tfrac12$; the atom bound is in
Lemma~\ref{lem:clt}.

(iii)  Fix $k$.  On the event $\{G(1)G(-1)\ne0\}$, the identity (i) gives
$\{\nu=k\}\subseteq\{\mathbf 1\{G(1)G(-1)<0\}\equiv k\bmod 2\}$, an event whose
probability tends to $\tfrac12$ regardless of the parity of $k$ by (ii).  Hence
\[
\PP(\nu=k)\;\le\;\tfrac12+o(1)+\PP\bigl(G(1)G(-1)=0\bigr)
\;=\;\tfrac12+o(1),
\]
uniformly in $k$.  For the stability statement: by
\eqref{eq:stability} and the winding-number formula, stability of
$(T_n(P_{r,s}))$ implies that $P_{r,s}$ has no zeros on $\T$ and
$\nu=r$; hence
$\PP(\text{stable})\le\PP(\nu=r)\le\tfrac12+o(1)$.
\end{proof}

\begin{lemma}[Reversal symmetry]\label{lem:reversal}
If the coefficient law is as in Theorem~\ref{thm:parity} and
$\PP(\text{$G_m$ has a zero on }\T)=0$ {\rm(}e.g.\ if $\xi_0$ has a
density{\rm)}, then $\nu$ and $m-\nu$ have the same distribution; in particular
$\E\nu=m/2$.
\end{lemma}

\begin{proof}
The reversed polynomial $G^*(z):=z^mG(1/z)=\sum_j\xi_{m-j}z^j$ has the same law
as $G$, and its zeros are the inverses of those of $G$; if no zeros lie on
$\T$, exactly $m-\nu$ of the inverted zeros lie in $\D$.
\end{proof}

\begin{corollary}[Band-reversal stability symmetry]\label{cor:band-reversal}
Fix a band split \(r+s=m\), and let the coefficients be i.i.d.\ from any
coefficient law.  Then the two complementary band sequences have the same
finite-degree stability probability:
\[
\PP\bigl((T_n(P_{r,s}))_{n\ge1}\ {\rm is\ stable}\bigr)
=
\PP\bigl((T_n(P_{s,r}))_{n\ge1}\ {\rm is\ stable}\bigr).
\]
If, in addition, \(\PP(\text{\(G_m\) has a zero on }\T)=0\), then
\[
\PP(\nu=k)=\PP(\nu=m-k)\qquad(0\le k\le m),
\]
and the common stability probability above is
\[
\PP(\nu=r)=\PP(\nu=s).
\]
\end{corollary}

\begin{proof}
For a deterministic coefficient vector put
\[
G^*(z):=z^mG(1/z),\qquad P^*_{s,r}(z):=z^{-s}G^*(z).
\]
Since \(m=r+s\),
\[
P^*_{s,r}(z)=z^{-s}z^mG(1/z)=z^rG(1/z)=P_{r,s}(1/z).
\]
Thus \(P^*_{s,r}\) has a zero on \(\T\) if and only if \(P_{r,s}\) does.  If
there is no zero on \(\T\), the path \(P^*_{s,r}|_{\T}\) is the reverse
parametrization of \(P_{r,s}|_{\T}\), so
\(\wind(P^*_{s,r})=-\wind(P_{r,s})\).  The stability criterion
\eqref{eq:stability} is therefore invariant under this reversal.

The reversed coefficient vector \((\xi_m,\xi_{m-1},\dots,\xi_0)\) has the same
law as the original i.i.d.\ vector, so the stability probability for the
reversed \((s,r)\)-symbol equals the stability probability for a fresh
\((s,r)\)-symbol.  This proves the first assertion.  Under the no-circle-zero
assumption, Lemma~\ref{lem:reversal} gives
\(\PP(\nu=k)=\PP(\nu=m-k)\).  The criterion \eqref{eq:stability} also makes the
\((r,s)\)-stability event equal to \(\{\nu=r\}\) almost surely, and the
\((s,r)\)-stability event equal in law to \(\{\nu=s\}\).
\end{proof}

\section{Strongly unbalanced bands: proof of Theorem \ref{thm:unbalanced}}
\label{sec:unbalanced}

\begin{proof}
Let $m=r+s$ and $|r-s|>2m^{9/10}$; without loss of generality
$r>\frac m2+m^{9/10}$.  Consider the three events
\[
E_1:=\Bigl\{\,\bigl|\nu-\tfrac m2\bigr|\le m^{9/10}\Bigr\},
\qquad
E_2:=\Bigl\{\,\text{no zeros of }G_m\text{ in }
\Bigl\{\bigl||z|-1\bigr|\le\tfrac{1}{m^{2}(\log m)^{3}}\Bigr\}\Bigr\}.
\]
By \cite[Theorem~3.1]{Manrique} (inequality \eqref{eq:window}),
$\PP(E_1^c)=O((\log m)^2m^{-1/10})$, and by the circle-clearance estimate
\cite[Theorem~1.1]{KMS}, as used in \cite[Section~2]{Manrique},
$\PP(E_2^c)=O((\log m)^{-1/2})$; hence
$\PP(E_1\cap E_2)=1-O((\log m)^{-1/2})$.

Work on $E_1\cap E_2$.  By $E_2$ the symbol $P_{r,s}(z)=z^{-r}G_m(z)$ has no
zeros on $\T$, so the winding number is defined and equals $\nu-r$.  By $E_1$
and the choice of $r$,
\[
\wind(P_{r,s})=\nu-r\;\le\;\tfrac m2+m^{9/10}-r\;<\;0 .
\]
Let $\delta:=\max\{|z|<1:G_m(z)=0\}$ and $\mu:=\min\{|z|>1:G_m(z)=0\}$ (with the
conventions $\delta=0$, $\mu=\infty$ if the corresponding set is empty; on
$E_1$ both sets are nonempty for large $m$ since $0<\nu<m$).  By $E_2$,
\[
\delta\le1-\frac{1}{m^2(\log m)^3},
\qquad
\mu\ge1+\frac{1}{m^2(\log m)^3},
\]
so
\[
\min\Bigl\{\log\frac1\delta,\ \log\mu\Bigr\}
\;\ge\;\log\Bigl(1+\frac{1}{m^{2}(\log m)^{3}}\Bigr)
\;>\;\alpha_m .
\]
By the exponential instability theorem for banded Toeplitz matrices with nonzero
winding number \cite[Theorem~4.1]{BG} (as quoted in
\cite[Section~1]{Manrique}), there is a constant $C=C(G_m)>0$ with
$\|T_n(P_{r,s})^{-1}\|\ge Ce^{\alpha_mn}$ for all $n\ge1$.  Finally, the
operator norm of a matrix dominates the modulus of each entry, so
$\sigma_1(T_n)\ge\max_{j}|\xi_j|$, which is positive on $E_1$ (there $G_m$ has
$m$ zeros, so $G_m\not\equiv0$); hence
$\kappa(T_n)=\sigma_1(T_n)\|T_n^{-1}\|\ge\bigl(\max_j|\xi_j|\bigr)Ce^{\alpha_mn}$
grows exponentially in $n$ on $E_1\cap E_2$.
\end{proof}

\section{The gap in the balanced case of \cite{Manrique}}\label{sec:gap}

We now explain why the balanced half of \cite[Theorem~2.1]{Manrique} is not
established by the argument given there, and why --- by
Theorem~\ref{thm:parity} --- no argument can establish it with
$c_3<\tfrac12$.

The proof of \cite[Theorem~2.1]{Manrique} begins by observing, correctly, that
$\wind(P)=0$ requires $r=\nu$.  It then invokes
\cite[Theorem~3.1]{Manrique} --- inequality \eqref{eq:window} above --- to
assert that with probability $1-O((\log m)^2m^{-1/10})$ the integer $r$
``might satisfy''
$\frac m2-m^{-1/10}\le r\le\frac m2+m^{-1/10}$, that the number of possible
integer values of $r$ is at most $2m^{-1/10}+1$, and hence that for large $m$
one has $r=\frac m2=\nu$.  But \eqref{eq:window} controls the
\emph{normalized} count $\nu/m$ to precision $m^{-1/10}$; for the count itself
this is the window
\[
\frac m2-m^{9/10}\;\le\;\nu\;\le\;\frac m2+m^{9/10},
\]
which contains $2m^{9/10}+1$ integers, not $2m^{-1/10}+1$.  Knowing that $\nu$
lies in a window of $\asymp m^{9/10}$ integers gives no lower bound on
$\PP(\nu=\frac m2)$: a priori the point probabilities could be as small as
$\asymp m^{-9/10}$.  The subsequent steps of the proof (the clearance of the
circle, the Wiener--Hopf bounds on $\|T(P)^{-1}\|$) are conditioned on the event
$\nu=r$ and are not in question; the gap is exactly the claim that this event
has probability close to $1$.

Theorem~\ref{thm:parity} shows the gap is essential: the event $\nu=\frac m2$
has probability at most $\tfrac12+o(1)$ for every coefficient law with zero
mean, unit variance and finite third moment, because the parity of $\nu$ is an
asymptotically fair coin.  Consequently the first display of
\cite[Theorem~2.1]{Manrique} --- positivity of $\lim_n\sigma_n(T_n)$, which by
\cite[Corollary~6.5]{BG} forces no zeros on $\T$ and zero winding, i.e.\
$\nu=r$ --- cannot hold with probability exceeding $\tfrac12+o(1)$, and the
suggestion of \cite[Remark~2.2]{Manrique} that the defect constant $c_3$ can be
made small (at the price of worsening other constants) is unfounded for
$c_3<\tfrac12$.

Two further remarks.  First, the parity obstruction is insensitive to the choice
of $r$: replacing the balanced band by $r=\frac m2\pm1$ merely flips the target
parity.  Second, the obstruction does not by itself decide whether
$\PP(\nu=\frac m2)$ is bounded below --- both
$\PP(\nu=\frac m2)\to\frac12$ and $\PP(\nu=\frac m2)\to0$ are consistent with
Theorem~\ref{thm:parity}; we expect the latter (Conjecture~\ref{conj:llt}).

\section{Towards the local limit theorem}\label{sec:open}

We close by recording what a resolution of Conjecture~\ref{conj:llt} would
require and the partial routes we see.

\begin{problem}\label{prob:variance}
Complete the conjectural linear variance asymptotics for
\(\sigma_m^2=\operatorname{Var}(\nu)\), where \(\nu\) is the number of zeros of a
Kac polynomial of degree \(m\) in the unit disk.
Proposition~\ref{prop:varupper} gives \(\sigma_m^2=O(m)\) for standard Gaussian
coefficients; prove the matching lower bound \(\sigma_m^2=\Omega(m)\), identify
the sharp constant \(c=\lim_m\sigma_m^2/m\), and decide whether central and local
limit theorems hold on that scale, here and for general coefficient laws.
\end{problem}

For real Gaussian coefficients the Kac--Rice machinery gives the zero-count
variance in closed integral form, and the upper bound of
Proposition~\ref{prop:varupper} can be read off from the resulting one- and
two-point intensities.

\begin{proof}[Proof of Proposition~\ref{prop:varupper}]
Write $N=\nu=N(\D)$ and let $\rho_1,\rho_2$ be the one- and two-point intensities
of the zero set of $G_m$ with respect to area measure $dA$.  The Kac--Rice second
moment identity for a point process is
\begin{equation}\label{eq:kr-var}
\operatorname{Var}(N)
=\int_{\D}\rho_1\,dA
+\iint_{\D\times\D}\bigl[\rho_2(z,w)-\rho_1(z)\rho_1(w)\bigr]\,dA(z)\,dA(w)
=:I_1+I_2,
\end{equation}
valid because $\E[N(N-1)]=\iint\rho_2$ and $\E N=\int\rho_1$; for Gaussian
$G_m$ the intensities are the Hammersley/Edelman--Kostlan permanental expressions
\cite{HKPV}.

\emph{Diagonal term.}  By the Edelman--Kostlan radial formula,
$\rho_1(z)=\frac1\pi\,\partial_z\partial_{\bar z}\log K_m(z,z)$ with
$K_m(z,w)=\sum_{j=0}^m(z\bar w)^j$, and integrating it over $\D$ gives the
classical count
\[
I_1=\int_\D\rho_1\,dA=\E N=\tfrac m2+O(\log m)=O(m).
\]

\emph{Connected term.}  Split $\D$ at radius $1-\frac{C\log m}{m}$.  On the bulk
$|z|,|w|\le1-\frac{C\log m}{m}$ the normalized kernel
$K_m(z,w)/\sqrt{K_m(z,z)K_m(w,w)}$ is uniformly close to the infinite-series
kernel $(1-z\bar w)^{-1}$ and the connected two-point function
$\rho_2-\rho_1\otimes\rho_1$ decays off the diagonal exactly as for the limiting
field, so its contribution to $I_2$ is $O(1)$.  What remains is the boundary
shell $S$ of radial width $\asymp1/m$, where $\rho_1\asymp m^2$.  The decisive
point is that on $S$ the connected kernel has \emph{no} $|z-w|^{-2}$ singularity:
by the local repulsion of complex Gaussian zeros \cite{NS},
$\rho_2(z,w)=O(|z-w|^2)$ as $w\to z$, so $\rho_2-\rho_1\otimes\rho_1$ stays
bounded on the diagonal rather than blowing up.  (This is the structural
difference from the one-dimensional real-axis Kac--Rice integral, whose
$|x-y|^{-2}$ near-diagonal behavior produces the $\log m$ of Maslova
\cite{Maslova}.)  A uniform $L^1$ decorrelation estimate then holds:
\[
\int_{\D}\bigl|\rho_2(z,w)-\rho_1(z)\rho_1(w)\bigr|\,dA(w)\;\le\;C\,\rho_1(z)
\qquad(z\in\D),
\]
with $C$ an absolute constant (numerically $\approx2$) independent of $m$,
obtained by integrating the off-diagonal exponential clustering of \cite{NS}
against the radial profile of $\rho_1$.  Hence
\[
|I_2|\le\int_\D\Bigl(\int_\D|\rho_2-\rho_1\otimes\rho_1|\,dA(w)\Bigr)dA(z)
\le C\int_\D\rho_1\,dA=C\,I_1=O(m).
\]
Combining the two bounds in \eqref{eq:kr-var} gives
$\operatorname{Var}(N)=O(m)$.
\end{proof}

The same identity points at the missing lower bound and the constant rather than
settling them.  After the edge rescaling, the near-circle zeros are modeled by the
Peres--Vir\'ag conformally invariant determinantal process
\cite{PV,HKPV}, whose Bergman kernel is
$\mathbb K(z,w)=\pi^{-1}(1-z\bar w)^{-2}$; for that process the disk count obeys
$\operatorname{Var}\bigl(N(|z|<r)\bigr)=r^2/(1-r^4)$ \cite{Buckley}, which at the
cutoff $r=1-1/m$ is $\approx m/4$.  This is a genuine cross-check of the order and
suggests a positive constant, but it is a statement about the limiting field, not
about the disk count of the finite Kac polynomial, so it does not by itself give
$\sigma_m^2=\Omega(m)$.  Two premise corrections are worth stating plainly.  The
$1/m$ above is the radial width of the shell that carries the $\Theta(m)$ zeros;
the spacing between consecutive zeros along a near-circle ray is finer, of order
$1/m^2$, and it is the shell width, not the spacing, that drives the variance.
And the proof above is for the \emph{real} Gaussian Kac polynomial used
throughout.  A genuinely complex-coefficient model has a different zero process,
without conjugate-pair symmetry, so constants suggested by that model should be
treated only as heuristic for the real disk count.  Numerically, direct simulation of
$\operatorname{Var}(\nu)/m$ over $m=25,\dots,400$ gives values in the range
$0.29$--$0.31$ with a log--log slope of $0.979$, consistent with
$\sigma_m^2\sim cm$, $c\approx0.30$.  Rigidity away from the boundary \cite{Yeager}
indicates annuli separated from $\T$ are not the source of large fluctuations, and
related boundary fluctuation constants in other ensembles are computed, for
example, in \cite{ACM}.  A natural route to the lower bound and to
anti-concentration is
through angular decompositions: the total winding of \(G_m\) along \(\T\) is the
sum of increments over \(K\) well-separated arcs, which should decorrelate because
the covariance kernel \(\sum_jz^j\bar w^j\) is a Dirichlet kernel, small for
angular separations \(\gg1/m\).  An Esseen-type estimate along this route could
give \(\max_k\PP(\nu=k)=O(K^{-1/2})\), ideally with \(K\) a positive power of
\(m\), or even \(K\asymp m\).  We hope to return to this elsewhere.

Finally, we note a statement of independent interest implied by
Theorem~\ref{thm:parity}(i)--(ii): \emph{for Kac polynomials with i.i.d.\
mean-zero, finite-third-moment coefficients, the number of zeros in the unit
disk is asymptotically equidistributed modulo $2$} --- in particular, for
Littlewood polynomials of degree $m$, the probability that exactly half of the
zeros lie inside the unit circle is at most $\tfrac12+o(1)$.

\subsection*{Acknowledgments}
The research presented in this paper was supported by the European Research Council
(ERC) under the European Union's Horizon Europe research and innovation programme
(grant agreement No.\ 101041711), by the Israel Science Foundation (grant number
2258/19) and by the Israel Science Foundation (ISF Grant 4101/25).

\subsection*{Declaration on generative AI}
During the preparation of this work the author used generative AI tools to
accelerate drafting and revision.  All mathematical claims, proofs, and
bibliographic information were subsequently reviewed and edited by the author,
who takes full responsibility for the content of the manuscript.

\begin{thebibliography}{9}

\bibitem{BG}
A.~B\"ottcher and S.~M.~Grudsky,
\emph{Spectral Properties of Banded Toeplitz Matrices},
SIAM, Philadelphia, 2005.

\bibitem{IZ}
I.~Ibragimov and O.~Zeitouni,
\emph{On roots of random polynomials},
Trans.\ Amer.\ Math.\ Soc.\ \textbf{349} (1997), no.~6, 2427--2441.

\bibitem{HN}
C.~P.~Hughes and A.~Nikeghbali,
\emph{The zeros of random polynomials cluster uniformly near the unit circle},
Compos.\ Math.\ \textbf{144} (2008), no.~3, 734--746.

\bibitem{KMS}
S.~V.~Konyagin and W.~Schlag,
\emph{Lower bounds for the absolute value of random polynomials on a
neighborhood of the unit circle},
Trans.\ Amer.\ Math.\ Soc.\ \textbf{351} (1999), no.~12, 4963--4980; we use the
strengthened form stated as Theorem~1.1 of reference [4] in \cite{Manrique}.

\bibitem{Manrique}
P.~Manrique,
\emph{The condition number of a random banded Toeplitz matrix is typically
large},
preprint, arXiv:2602.12581 (2026).

\bibitem{Shevtsova}
I.~Shevtsova,
\emph{On the absolute constants in the Berry--Esseen-type inequalities},
Dokl.\ Math.\ \textbf{89} (2014), 378--381.

\bibitem{Maslova}
N.~B.~Maslova,
\emph{The variance of the number of real roots of random polynomials},
Theory Probab.\ Appl.\ \textbf{19} (1974), no.~1, 35--52.

\bibitem{DoNguyen}
Y.~Q.~Do and N.~D.~V.~Nguyen,
\emph{Real roots of random polynomials: asymptotics of the variance},
Electron.\ J.\ Probab.\ \textbf{30} (2025), paper no.~112; arXiv:2303.05478.

\bibitem{Yeager}
A.~M.~Yeager,
\emph{The variance of the number of zeros for complex random polynomials
spanned by OPUC},
preprint, arXiv:1908.02234 (2019).

\bibitem{HKPV}
J.~B.~Hough, M.~Krishnapur, Y.~Peres and B.~Vir\'ag,
\emph{Zeros of Gaussian Analytic Functions and Determinantal Point Processes},
University Lecture Series, vol.~51, Amer.\ Math.\ Soc., Providence, RI, 2009.

\bibitem{NS}
F.~Nazarov and M.~Sodin,
\emph{Correlation functions for random complex zeroes: strong clustering and
local universality},
Comm.\ Math.\ Phys.\ \textbf{310} (2012), no.~1, 75--98.

\bibitem{PV}
Y.~Peres and B.~Vir\'ag,
\emph{Zeros of the i.i.d.\ Gaussian power series: a conformally invariant
determinantal process},
Acta Math.\ \textbf{194} (2005), no.~1, 1--35.

\bibitem{Buckley}
J.~Buckley,
\emph{Fluctuations in the zero set of the hyperbolic Gaussian analytic
function},
Int.\ Math.\ Res.\ Not.\ IMRN \textbf{2015}, no.~6, 1666--1687.

\bibitem{SV}
L.~A.~Shepp and R.~J.~Vanderbei,
\emph{The complex zeros of random polynomials},
Trans.\ Amer.\ Math.\ Soc.\ \textbf{347} (1995), no.~11, 4365--4384.

\bibitem{ACM}
Y.~Ameur, C.~Charlier and P.~Moreillon,
\emph{Eigenvalues of truncated unitary matrices: disk counting statistics},
preprint, arXiv:2305.08976 (2023).

\bibitem{condrefs}
See the references [3,~16,~22,~23] of \cite{Manrique} for the conditioning of
dense random Toeplitz and circulant matrices.

\end{thebibliography}

\end{document}
